\documentclass[
  lang = cn,
  theme = Fresh Green,
  font = ctex-fandol,
  mode = fancy,
  thmcnt = chapter,
  gb7714-2015,
  bibend = biber
]{elegantbook-l3}

\usepackage{amsmath, amssymb, array}
\addbibresource{biblatex-examples.bib}

\elegantnewtheorem{sampletheorem}{自定义定理}{thmstyle}{sample}
\elegantnewtheorem{sampledefinition}{自定义定义}{defstyle}{sampledef}

\begin{document}

\ElegantCoverInfo{%
  title = {ElegantBook LaTeX3 迁移验证},
  subtitle = {当前已实现功能的中文回归测试文档},
  cover image = {cover.jpg},
  logo image = {logo-blue},
  author = {ElegantBook LaTeX3 Refactor},
  institute = {ElegantBook Program},
  version = {v1.0},
  date = {2026/05/04},
  extrainfo = {本文件用于验证当前重构版已经迁移的功能。}
}

\maketitle
\frontmatter
\tableofcontents

\mainmatter

\chapter{结构与基础排版}

\section{层级标题}

本节验证章、节、小节和小小节标题的编号、颜色、页眉标记与目录条目。

\subsection{小节层级}

正文段落用于观察基础行距、缩进、中文字体和西文字体的协调性。当前主题色应统一作用于章节标题、页眉线、链接以及后续环境标题。

\subsubsection{小小节层级}

数学公式用于验证基础数学包、编号和交叉引用的配合。式 \eqref{eq:quadratic} 给出一个标准恒等式：
\begin{equation}
  (a+b)^2 = a^2 + 2ab + b^2 .
  \label{eq:quadratic}
\end{equation}

\section{列表与链接}

列表项目符号应呈现主题色，嵌套层级应保持稳定缩进。

\begin{itemize}
  \item 一级无序列表项目。
  \begin{itemize}
    \item 二级无序列表项目。
    \begin{itemize}
      \item 三级无序列表项目。
    \end{itemize}
  \end{itemize}
\end{itemize}

\begin{enumerate}
  \item 一级有序列表项目。
  \item 第二个有序列表项目，包含行内链接：
    \href{https://github.com/ElegantLaTeX/ElegantBook}{ElegantBook Repository}。
\end{enumerate}

\section{章节引言}

\begin{introduction}
  \item 引言环境应使用主题色标题和项目符号。
  \item 该环境应能容纳多条摘要式内容。
  \item 盒子宽度应与正文区域对齐。
\end{introduction}

\chapter{定理与普通环境}

\section{盒式定理类环境}

\begin{definition}[度量空间]\label{def:metric-space}
  设 $X$ 为集合，若函数 $d:X\times X\to \mathbb{R}$ 满足非负性、同一性、对称性和三角不等式，则称 $(X,d)$ 为度量空间。
\end{definition}

\begin{theorem}{闭区间连续函数有界}{bounded-continuous}
  若函数 $f$ 在闭区间 $[a,b]$ 上连续，则 $f$ 在 $[a,b]$ 上有界。
\end{theorem}

\begin{lemma}[三角不等式的直接推论]\label{lem:triangle}
  对任意实数 $x,y$，有 $|x+y|\le |x|+|y|$。
\end{lemma}

\begin{corollary}[连续函数局部有界]
  若函数在某点连续，则它在该点的某个邻域内有界。
\end{corollary}

\begin{postulate}[平行公设]
  过直线外一点有且只有一条直线与已知直线平行。
\end{postulate}

\begin{axiom}[选择公理]
  任意一族非空集合存在选择函数。
\end{axiom}

\begin{proposition}[单调有界收敛]
  单调有界数列必收敛。
\end{proposition}

引用检查：定义 \ref{def:metric-space}、定理 \ref{thm:bounded-continuous}、引理 \ref{lem:triangle}。

\begin{theorem*}{无编号定理}
  该环境验证星号版本不参与编号，但保持相同的视觉风格。
\end{theorem*}

\section{导言区声明的定理环境}

\begin{sampletheorem}{自定义旧式标题}{custom-theorem}
  该环境由导言区的 \lstinline|\elegantnewtheorem| 声明，采用旧式标题和标签参数。
\end{sampletheorem}

\begin{sampledefinition}[自定义可选标题]\label{sampledef:native}
  该环境验证自定义定义类环境的可选标题与原生标签写法。
\end{sampledefinition}

\begin{sampletheorem*}[自定义无编号环境]
  该环境验证自定义定理的星号版本。
\end{sampletheorem*}

自定义环境引用检查：定理 \ref{sample:custom-theorem}、定义 \ref{sampledef:native}。

\section{非盒式普通环境}

\begin{example}[紧性示例]\label{exam:compact}
  闭区间 $[0,1]$ 是实数通常拓扑下的紧集。
\end{example}

\begin{exercise}[基础计算]\label{exer:basic}
  计算 $\int_0^1 x^2\,dx$。
\end{exercise}

\begin{problem}[极值问题]\label{prob:extreme}
  求函数 $f(x)=x^2-2x+1$ 在实数域上的最小值。
\end{problem}

普通编号环境引用检查：例题 \ref{exam:compact}、练习 \ref{exer:basic}、问题 \ref{prob:extreme}。

\begin{note}
  笔记环境应以左侧符号和粗体引导词开始，标题文字与其他普通环境左边界对齐。
\end{note}

\begin{proof}
  证明环境应以粗体引导词开始，并使用当前设定的正文字体风格。
\end{proof}

\begin{solution}
  解答环境用于呈现答案内容；当前文档验证其默认显示状态。
\end{solution}

\begin{remark}
  注环境是普通段落式环境，不应被定义为盒式定理环境。
\end{remark}

\begin{assumption}
  假设环境用于列出后续论证所依赖的条件。
\end{assumption}

\begin{conclusion}
  结论环境用于给出阶段性推论。
\end{conclusion}

\begin{property}
  性质环境用于描述对象的稳定属性。
\end{property}

\begin{custom}{自定义普通环境}
  自定义普通环境应接受一个标题参数，并保持与结论类环境一致的段落样式。
\end{custom}

\begin{problemset}[本章综合练习]
  \item 验证闭区间上连续函数有界。
  \item 构造一个有界但不闭的实数子集。
  \item 比较例题、练习和问题三个计数器的显示格式。
\end{problemset}

\chapter{代码高亮与跨页}

\section{行内与短代码块}

行内代码用于验证 \lstinline|\lstinline| 的字体、断行和颜色设置。下面的短代码块验证默认 LaTeX 语言高亮：

\begin{lstlisting}
\documentclass[lang=cn, theme=Fresh Green]{elegantbook-l3}
\begin{document}
\chapter{Minimal Test}
\end{document}
\end{lstlisting}

命令行代码块用于验证显式语言选项：

\begin{lstlisting}[language=bash]
latexmk -pdfxe elegantbook-cn.tex
\end{lstlisting}

\section{跨页代码块}

下面的长代码块用于验证普通 \lstinline|lstlisting| 能够自然跨页。

\begin{lstlisting}
\ExplSyntaxOn
\cs_new_protected:Npn \demo_listing_cross_page_line:n #1
  {
    \iow_term:n { Listing cross page test line #1 }
  }
\demo_listing_cross_page_line:n { 001 }
\demo_listing_cross_page_line:n { 002 }
\demo_listing_cross_page_line:n { 003 }
\demo_listing_cross_page_line:n { 004 }
\demo_listing_cross_page_line:n { 005 }
\demo_listing_cross_page_line:n { 006 }
\demo_listing_cross_page_line:n { 007 }
\demo_listing_cross_page_line:n { 008 }
\demo_listing_cross_page_line:n { 009 }
\demo_listing_cross_page_line:n { 010 }
\demo_listing_cross_page_line:n { 011 }
\demo_listing_cross_page_line:n { 012 }
\demo_listing_cross_page_line:n { 013 }
\demo_listing_cross_page_line:n { 014 }
\demo_listing_cross_page_line:n { 015 }
\demo_listing_cross_page_line:n { 016 }
\demo_listing_cross_page_line:n { 017 }
\demo_listing_cross_page_line:n { 018 }
\demo_listing_cross_page_line:n { 019 }
\demo_listing_cross_page_line:n { 020 }
\demo_listing_cross_page_line:n { 021 }
\demo_listing_cross_page_line:n { 022 }
\demo_listing_cross_page_line:n { 023 }
\demo_listing_cross_page_line:n { 024 }
\demo_listing_cross_page_line:n { 025 }
\demo_listing_cross_page_line:n { 026 }
\demo_listing_cross_page_line:n { 027 }
\demo_listing_cross_page_line:n { 028 }
\demo_listing_cross_page_line:n { 029 }
\demo_listing_cross_page_line:n { 030 }
\demo_listing_cross_page_line:n { 031 }
\demo_listing_cross_page_line:n { 032 }
\demo_listing_cross_page_line:n { 033 }
\demo_listing_cross_page_line:n { 034 }
\demo_listing_cross_page_line:n { 035 }
\demo_listing_cross_page_line:n { 036 }
\demo_listing_cross_page_line:n { 037 }
\demo_listing_cross_page_line:n { 038 }
\demo_listing_cross_page_line:n { 039 }
\demo_listing_cross_page_line:n { 040 }
\demo_listing_cross_page_line:n { 041 }
\demo_listing_cross_page_line:n { 042 }
\demo_listing_cross_page_line:n { 043 }
\demo_listing_cross_page_line:n { 044 }
\demo_listing_cross_page_line:n { 045 }
\demo_listing_cross_page_line:n { 046 }
\demo_listing_cross_page_line:n { 047 }
\demo_listing_cross_page_line:n { 048 }
\demo_listing_cross_page_line:n { 049 }
\demo_listing_cross_page_line:n { 050 }
\demo_listing_cross_page_line:n { 051 }
\demo_listing_cross_page_line:n { 052 }
\demo_listing_cross_page_line:n { 053 }
\demo_listing_cross_page_line:n { 054 }
\demo_listing_cross_page_line:n { 055 }
\demo_listing_cross_page_line:n { 056 }
\demo_listing_cross_page_line:n { 057 }
\demo_listing_cross_page_line:n { 058 }
\demo_listing_cross_page_line:n { 059 }
\demo_listing_cross_page_line:n { 060 }
\demo_listing_cross_page_line:n { 061 }
\demo_listing_cross_page_line:n { 062 }
\demo_listing_cross_page_line:n { 063 }
\demo_listing_cross_page_line:n { 064 }
\demo_listing_cross_page_line:n { 065 }
\demo_listing_cross_page_line:n { 066 }
\demo_listing_cross_page_line:n { 067 }
\demo_listing_cross_page_line:n { 068 }
\demo_listing_cross_page_line:n { 069 }
\demo_listing_cross_page_line:n { 070 }
\demo_listing_cross_page_line:n { 071 }
\demo_listing_cross_page_line:n { 072 }
\demo_listing_cross_page_line:n { 073 }
\demo_listing_cross_page_line:n { 074 }
\demo_listing_cross_page_line:n { 075 }
\demo_listing_cross_page_line:n { 076 }
\demo_listing_cross_page_line:n { 077 }
\demo_listing_cross_page_line:n { 078 }
\demo_listing_cross_page_line:n { 079 }
\demo_listing_cross_page_line:n { 080 }
\ExplSyntaxOff
\end{lstlisting}

\chapter{表格与颜色}

\section{主题颜色}

下表用于验证当前主题导出的兼容颜色名是否可用于普通文本、表格和链接上下文。

\begin{center}
\begin{tabular}{>{\bfseries}m{0.24\textwidth}m{0.55\textwidth}}
\hline
颜色名 & 验证文本 \\
\hline
structurecolor & \textcolor{structurecolor}{章节结构色与页眉线颜色一致。} \\
main & \textcolor{main}{定义、例题、练习与问题使用该主色。} \\
second & \textcolor{second}{定理、笔记、证明与注环境使用该辅助色。} \\
third & \textcolor{third}{命题、假设、结论、性质与自定义普通环境使用该颜色。} \\
winered & \textcolor{winered}{代码高亮中的控制序列和关键字颜色。} \\
\hline
\end{tabular}
\end{center}

\section{排版收束}

本文档只覆盖当前已经迁移并用于回归检查的功能。尚未迁移的版本历史、完整目录附录样式和旧封面命令不会在此文件中作为测试目标。

\chapter{参考文献}

\section{引用与文献表}

本节使用 TeX Live 自带的 \lstinline|biblatex-examples.bib| 验证参考文献接口。行文引用示例包括 Knuth 的多卷本 \cite{knuth:ct:a,knuth:ct:b}、\TeX{} 参考书 \cite{companion}，以及一篇期刊论文 \cite{sigfridsson}。

也可以在同一句中混合多条文献，观察通过类选项调用的 GB/T 7714-2015 引用样式是否生效 \cite{aksin,angenendt,bertram}.

\printbibliography

\end{document}
